% ===================================================================
%  2. TAULA — Giza gorputza. — Jolasaldia. — Jokoak.
%  Traduction 2026 d'apres texto/io, controlee sur texto/fr, texto/en
%  et texto/es. Consigne : voir texto/eu/10-taula-01.tex.
%
%  LE TABLEAU DE L'ANATOMIE EST CELUI OU L'ON MESURE QUE LE BASQUE NE
%  DOIT RIEN A PERSONNE. Pas un nom du corps ne rappelle un mot connu :
%  la tete est « burua », le bras « besoa », la jambe « hanka », le
%  cœur « bihotza », l'œil « begia », la main « eskua ». Quinze
%  colonnes servies, et aucune n'aide a lire celle-ci.
%
%  UN MOT DE LA MAISON, POURTANT. Le jeu du bloc 24 — « pilota hormaren
%  kontra », la balle contre le mur — est justement la pelote basque,
%  et c'est le seul endroit du volume ou le jeu decrit par Guignon
%  porte, dans cette colonne, son nom propre plutot qu'une traduction.
% ===================================================================

%%K t02-apar-1 apar
\VUsaut{4.0mm}
\VUtitre{15.0pt}{60}{2. TAULA}
\VUsaut{2.0mm}
\VUtitre{11.4pt}{0}{Giza gorputza. --- Jolasaldia.}
\VUtitre{11.4pt}{0}{Jokoak.}

%%K t02-tit-0 sub
\VUtitre{13.2pt}{0}{Giza gorputza.}
\VUsaut{2.4mm}
\VUcentre{10.2pt}{0}{\textit{(Natur zientzien ikasgai erraza.)}}
\VUsaut{3.0mm}

%%K t02-01-1 p
1. --- Giza gorputzaren zati nagusiak hauek dira: \VUgras{burua},
\VUgras{enborra} eta \VUgras{gorputz-adarrak}.

%%K t02-tit-1 sub
\VUpk{10.2pt}{40}{Burua.}
\VUsaut{3.0mm}

%%K t02-02-1 p
2. --- Buruak bi zati ditu: \VUgras{garezurra}, zeinaren barruan \VUgras{garuna}
dagoen eta zeina \VUgras{ileak} estaltzen duen, eta \VUgras{aurpegia}.

%%K t02-03-1 p
3. --- Gure marrazkian ez da garezurrik ez garunik ikusten; aitzitik, oso argi
ikusten dira aurpegiaren zati garrantzitsuak.

%%K t02-04-1 p
4. --- Hona hemen lehenik \VUgras{bekokia}, adimen sendoa eta etengabe lanean ari
den pentsamendua adierazten dituena. \VUgras{Lokiak} oso agerian daude.
\VUgras{Begia} bizia eta zabal irekia da. \VUgras{Bekain} lodi batek eta
\VUgras{betile} luzeek babesten dute.

%%K t02-05-1 p
5. --- Hona hemen \VUgras{sudurra} eta \VUgras{sudur-zuloak} ere. Sudurra
usaintzeko eta arnasa hartzeko balio digu. Sudurraren bi aldeetan
\VUgras{masailak} daude, eta urrunago \VUgras{belarriak}. Aurpegiaren beheko
zatia \VUgras{ahoak} eta \VUgras{kokotsak} osatzen dute.

%%K t02-06-1 p
6. --- Gaizki ikusten dira, \VUgras{biboteak} eta \VUgras{bizarrak} ezkutatzen
baitizkigute. Baina badakigu ahoan bi \VUgras{ezpain} daudela, \VUgras{goiko
masailezurra} eta \VUgras{beheko masailezurra} beren \VUgras{hortzekin}, eta
\VUgras{mihia}. Ahoaz hitz egiten eta jaten dugu. Mihiaz janariaren zaporea
sentitzen dugu.

%%K t02-07-1 p
7. --- Begia, belarria, sudurra eta ahoa, hatzekin batera, bost zentzumenen
organoak dira: \VUgras{ikusmenarena}, \VUgras{entzumenarena},
\VUgras{usaimenarena}, \VUgras{dastamenarena} eta \VUgras{ukimenarena}.

%%K t02-08-1 p
8. --- Burua da, beraz, giza gorputzaren zatirik bitxienetako bat.

%%K t02-tit-2 sub
\VUsaut{4.4mm}
\VUpk{10.2pt}{40}{Enborra.}
\VUsaut{3.0mm}

%%K t02-09-1 p
9. --- \VUgras{Lepoak} burua enborrarekin lotzen du. Haren barnean
\VUgras{eztarria} eta \VUgras{garondoa} daude.

%%K t02-10-1 p
10. --- Enborrean ere organo beharrezko asko dago. \VUgras{Bularrean}
\VUgras{birikak} daude, zeinekin arnasa hartzen dugun, eta \VUgras{bihotza},
bizitzaren erdigunea. Bihotzak taupadak eteten dituenean, izaki bizidunak hil
egiten dira.

%%K t02-11-1 p
11. --- \VUgras{Saihetsek} bularra eusten dute.

%%K t02-12-1 p
12. --- Enborraren beste zatiak \VUgras{saihetsaldea}, \VUgras{bizkarra},
\VUgras{sorbaldak} eta \VUgras{sabela} dira. Saihetsaldean edo bizkarrean etzaten
gara lo egiteko, eta zamak sorbaldan eramaten ditugu.

%%K t02-13-1 p
13. --- Sabelean \VUgras{urdaila} eta \VUgras{hesteak} daude. Janaria digeritzeko
balio digute.

%%K t02-tit-3 sub
\VUsaut{4.4mm}
\VUpk{10.2pt}{40}{Gorputz-adarrak.}
\VUsaut{3.0mm}

%%K t02-14-1 p
14. --- Lau \VUgras{gorputz-adar} ditugu: bi \VUgras{beso} eta bi \VUgras{hanka}.
Besoekin gauzak hartzen, eusten, jasotzen eta biltzen ditugu; hankekin zutik
egoten, ibiltzen eta korrika egiten dugu.

%%K t02-15-1 p
15. --- \VUgras{Sorbaldak} besoa gorputzarekin lotzen du. \VUgras{Gihar} oso
sendo batek, biburuduna, besoa mugitzen du, eta \VUgras{ukondoak} besoa
\VUgras{besaurrearekin} lotzen du. \VUgras{Eskumuturrak} \VUgras{eskuaren}
artikulazioa osatzen du, eta eskua \VUgras{hatzetan} eta \VUgras{azazkaletan}
amaitzen da. Eskuan \VUgras{zaina} ikusten da (eta arteria), zeinetik
\VUgras{odola} igarotzen den.

%%K t02-16-1 p
16. --- Hankaren zatiak hauek dira: \VUgras{izterra}, \VUgras{belauna} eta
\VUgras{belaun-zuloa}, \VUgras{bernazakia}, zeinaren \VUgras{haragia}
\VUgras{larruazalak} estaltzen duen, \VUgras{orkatila} eta \VUgras{oina}.
Hankaren \VUgras{hezurrak} oso lodiak eta oso sendoak dira. Oinaren muturrean
\VUgras{oin-hatzak} daude. Beste zatiak \VUgras{zola} eta \VUgras{orpoa} dira.

%%K t02-17-1 p
17. --- Halakoa da giza gorputzaren deskribapen ia osoa.

%%K t02-17-2 p
\textit{Oharra.} --- \textit{«(burua eta abar) min dut» esapidearen bidez
irakasleak hiztegi hau guztia berriro landu ahal izango du} \VUgras{gorputzeko
osasunaren} \textit{aldetik, ona edo txarra}.

%%K t02-tit-4 sub
\VUsaut{5.0mm}
\VUpk{10.2pt}{40}{Gimnasia.}
\VUsaut{2.4mm}
\VUcentre{10.2pt}{0}{\textit{(Ikasleetako batek kontatua.)}}
\VUsaut{3.0mm}

%%K t02-18-1 p
18. --- Malguak, trebeak eta osasuntsuak izateko, hankak eta besoak landu behar
dira. Horregatik jolasten gara eta \VUgras{gimnasia} egiten dugu.

%%K t02-19-1 p
19. --- Astean zehar bi jarduera horiek eskolako jolastokian egiten dira (ikus 1.
taula). Baina ostegunetan eta igandeetan zelai handi batera ateratzen gara, non
korrika egiteko eta ondo pasatzeko lekua baitago.

%%K t02-20-1 p
20. --- Gure irakasleak eta gurasoak batzuetan gurekin etortzen dira; baina
ariketak \VUgras{gimnasia-irakasleak} \textsuperscript{(12)} zuzentzen ditu.

%%K t02-21-1 p
21. --- Zelaian, sarreraren ezkerrean, \VUgras{gimnasia-tresnak}
\textsuperscript{(1)} daude: \VUgras{eskaileratik} \textsuperscript{(2)} gora
igotzen gara, \VUgras{eraztunetan} \textsuperscript{(3)} eta \VUgras{trapezioan}
\textsuperscript{(4)} burua behera zabuka aritzen gara, \VUgras{soka-eskaileraren}
\textsuperscript{(5)} edo \VUgras{korapilodun sokaren} \textsuperscript{(6)}
gaineraino besoez igotzen gara.

%%K t02-22-1 p
22. --- Bitartean gure lagunak \VUgras{zurezko zaldiaren} \textsuperscript{(7)}
edo \VUgras{barren} \textsuperscript{(11)} gainetik jauzi egiten dute. Beste
batzuk \VUgras{boxeatzen} \textsuperscript{(8)} edo \VUgras{eskriman}
\textsuperscript{(9)} aritzen dira, maskara jantzita eta \VUgras{floretarekin}.
Beste batzuek, azkenik, \VUgras{halterak} \textsuperscript{(10)} jasotzen
dituzte.

%%K t02-tit-5 sub
\VUsaut{4.6mm}
\VUpk{10.2pt}{40}{Jokoak.}
\VUsaut{3.0mm}

%%K t02-23-1 p
23. --- Gimnasten inguruan mutilak eta neskak \VUgras{joko} askotara jolasten
dira. Neskak beren \VUgras{uztaiarekin} \textsuperscript{(14)} olgatzen dira;
\VUgras{soka-saltoan} \textsuperscript{(13)} aritzen dira, edo nebei eta
lehengusinei deitzen diete \VUgras{kroket} \textsuperscript{(15)} partida bat
egitera. Zein pozik jotzen duten aurkariaren \VUgras{bola} beren \VUgras{zurezko
mailuaz}, hark arkua erraz igaroko duela uste duen unean!

%%K t02-24-1 p
24. --- Mutilek joko bizkorragoak nahiago dituzte. Gogotik jolasten dira
\VUgras{boloetara} \textsuperscript{(16)}, zeinak trebeki botatzen baitituzte
\VUgras{bolaz} \textsuperscript{(17)}. Ez dituzte gutxiesten \VUgras{ziba}
\textsuperscript{(18)} eta \VUgras{bilboketa} \textsuperscript{(19)} ere, ezta
\VUgras{igela} \textsuperscript{(20)} ere. Baina beren jokorik maiteenak
\VUgras{kanikak} \textsuperscript{(21)} eta \VUgras{pilota hormaren kontra}
\textsuperscript{(22)} dira, \VUgras{negutegiaren} \textsuperscript{(26)} horma
egokia baita pilota arina harengandik itzultzeko.

%%K t02-25-1 p
25. --- Joanes txikia oso trebea da bere \VUgras{makil-hanketan}
\textsuperscript{(24)}, eta oreka bikain eusten du. Haren ondoan lagun batzuk
\VUgras{ezkutaketan} \textsuperscript{(25)} ari dira negutegi horretan eta
\VUgras{hesiaren} atzean. Beste batzuek \VUgras{itsumandoka eskuz}
\textsuperscript{(28)} edo \VUgras{bolantea} \textsuperscript{(29)} nahiago
dituzte, zeina \VUgras{erraketa} batetik bestera hegan baitoa.

%%K t02-noto-1 noto
\VUnotes{143.2mm}{%
(*) Haurren jokoek maiz izen erabat nazionalak dituzte. Idoz izendatu ditugu
ahal izan dugun bezala: batzuetan bereizten dituen ekintzagatik, beste
batzuetan herrialde baten edo bestearen izena hartuz, nahasgarriegia ez zenean.
Adibidez: \VUgras{upela} (alemana eta frantsesa) \VUgras{igela}
\textsuperscript{(20)} da (ingelesa), non jokalariek beren txanpon biribilak
metalezko igel baten ahoan sartzen saiatzen baitira, zulodun kutxa (upel) baten
gainean zabalik. \VUgras{Sorbalda gaineko jauzia} \textsuperscript{(30)}
\VUgras{astokatik} bereizten da soilik honengatik: buruaren gainetik jauzi
egitean, jokalariek eskuak lagunaren sorbaldetan bermatzen dituzte,
«\VUgras{astokan}» aldiz bizkarrean bakarrik bermatzen dituzten bitartean. ---
Edo bestela: jokalari batzuk makurtzen dira, eta besteak beren bizkarren
gainetik jauzi egiten dute, eskuak sorbaldetan bermatuz; lehenengoek bultzada
jasan behar dute mugitu gabe, bigarrenek oreka galdu gabe jauzi egin (ikus taula
bera).%
}

%%K t02-26-1 p
26. --- Zelaiaren erdian bi zuhaitz hazten dira. Haietako baten oinean zazpi edo
zortzi mutilek \VUgras{sorbalda gaineko jauzian} \textsuperscript{(30)} (*) hasi
dute jokoa. Joko hau ez da herrialde guztietan ezaguna; baina denek dakite zer
den \VUgras{astoka}, zeinera haurrak orain ez baitira jolasten.

%%K t02-tit-6 sub
\VUsaut{5.0mm}
\VUpk{10.2pt}{40}{Aire zabaleko jokoak.}
\VUsaut{3.4mm}

%%K t02-27-1 p
27. --- \VUgras{Itsumandoka} \textsuperscript{(31)} ere joko oso alaia da; neskak
eta mutilak txandaka \VUgras{zapia} jartzen dute begietan eta harrapatzen uzten
duten baldarrak harrapatzen dituzte.

%%K t02-28-1 p
28. --- Zelaiaren erdian --- hona hemen joko erabat frantsesa, zakar samarra
bada ere: \VUgras{hartza} \textsuperscript{(32)} da, \VUgras{kometaren}
\textsuperscript{(33)} joko baketsua baino askoz zaratatsuagoa, eta ingelesen
\VUgras{tenis} \textsuperscript{(34)} jokoa baino ere bai.

%%K t02-29-1 p
29. --- Azken kirol hau ikasle handienen atsegina da. Gure irakasleek berek ere
gogotik heltzen diote. Trebetasun eta bizkortasun handia eskatzen du, eta oso
maite dut. \VUgras{Zabua} \textsuperscript{(35)} neskei uzten diet.

%%K t02-30-1 p
30. --- Ez zait gustatzen arriskutsu bihur daitekeen beste ingeles joko bat ere.

%%K t02-30-1b p
\VUgras{Futbolaz} \textsuperscript{(36)} ari naiz, nire anaia zaharrenaren
poza. Nahiago dut gure \VUgras{presoen} \textsuperscript{(38)} joko zaharra,
osasuntsu bezain zakartasun gutxiagokoa.

%%K t02-tit-7 sub
\VUsaut{4.6mm}
\VUpk{10.2pt}{40}{Bizikleta eta pilota-jokoak.}
\VUsaut{3.0mm}

%%K t02-31-1 p
31. --- Neuk ere gogotik ematen diot itzulia zelaiari \VUgras{bizikletan}
\textsuperscript{(37)}.

%%K t02-32-1 p
32. --- Datorren urtean \VUgras{zaldiz ibiltzen} \textsuperscript{(39)} ikasiko
dut. Zein ederra den zaldia, pausoan edo trostan dotore doanean eta oztopoak
lauhazka gainditzen dituenean!

%%K t02-33-1 p
33. --- Udan \VUgras{igeri egiten} \textsuperscript{(40)} ikasten dugu. Urpean
sartzen gara eta gozamenez bainatzen ibaiko ur garbi eta freskoan. Batzuetan
amaieran paperezko \VUgras{globo} bat \textsuperscript{(41)} botatzen dugu, zeina
handikiro igotzen baita airera, aire beroz bete bezain laster.

%%K t02-34-1 p
34. --- Beste joko asko ere badago, baina inoiz ez nuke amaituko haiek
zerrendatzen, eta kontakizun labur honetatik ageri da gure zelai ederrean
ostegunak batere aspergarriak ez direla.

%%K t02-34-2 p
N.~B. --- \textit{Irakasleari ez zaio zaila izango kapitulu hau osatzea, jokorik
garrantzitsuenetako bakoitza zehatzago deskribatuz.}
\VUsaut{6.0mm}
\VUfilet{22mm}
